% Synthetic fixture for the article-structure sidecars (Spectra change
% add-article-construction-model, task 1.3). Invented text on a memory-research
% theme; it is not taken from any manuscript. One proposition per line.
\documentclass{article}
\title{Spaced Review and Delayed Recall}
\begin{document}
\maketitle

\section{Introduction}\label{sec:introduction}
Spaced review is widely recommended, but earlier comparisons often gave the spaced group more total study time than the massed group.
We therefore ask whether spacing improves recall one week later when total study time is equal.
In other words, the main goal of this study is to compare spaced and massed review at equal study time.
We expected spaced review to produce higher recall than massed review,
because each spaced review requires retrieval after a delay.
The rest of this article describes the method, reports the results, and closes with a discussion.

\section{Method}\label{sec:method}
Twenty-four adult volunteers reviewed forty word pairs, either three times in one session or once on each of three days.

\section{Results}\label{sec:results}
One week later, mean recall was 0.62 after spaced review and 0.41 after massed review.
These results support the prediction that spacing improves delayed recall.

\subsection{An alternative explanation}\label{sec:alternative}
However, longer total study time could also produce such an advantage, an alternative we now examine.

\section{Discussion}\label{sec:discussion}
As noted in the Introduction, earlier comparisons confounded spacing with total study time, but here both groups studied for the same total time, so study time cannot explain the advantage.

In addition, future studies could test whether the advantage lasts beyond one week.

\end{document}
